% Options for packages loaded elsewhere
\PassOptionsToPackage{unicode}{hyperref}
\PassOptionsToPackage{hyphens}{url}
\PassOptionsToPackage{dvipsnames,svgnames,x11names}{xcolor}
%
\documentclass[
  letterpaper,
  twocolumn,
  landscape]{report}

\usepackage{amsmath,amssymb}
\usepackage{iftex}
\ifPDFTeX
  \usepackage[T1]{fontenc}
  \usepackage[utf8]{inputenc}
  \usepackage{textcomp} % provide euro and other symbols
\else % if luatex or xetex
  \usepackage{unicode-math}
  \defaultfontfeatures{Scale=MatchLowercase}
  \defaultfontfeatures[\rmfamily]{Ligatures=TeX,Scale=1}
\fi
\usepackage{lmodern}
\ifPDFTeX\else  
    % xetex/luatex font selection
\fi
% Use upquote if available, for straight quotes in verbatim environments
\IfFileExists{upquote.sty}{\usepackage{upquote}}{}
\IfFileExists{microtype.sty}{% use microtype if available
  \usepackage[]{microtype}
  \UseMicrotypeSet[protrusion]{basicmath} % disable protrusion for tt fonts
}{}
\makeatletter
\@ifundefined{KOMAClassName}{% if non-KOMA class
  \IfFileExists{parskip.sty}{%
    \usepackage{parskip}
  }{% else
    \setlength{\parindent}{0pt}
    \setlength{\parskip}{6pt plus 2pt minus 1pt}}
}{% if KOMA class
  \KOMAoptions{parskip=half}}
\makeatother
\usepackage{xcolor}
\setlength{\emergencystretch}{3em} % prevent overfull lines
\setcounter{secnumdepth}{5}
% Make \paragraph and \subparagraph free-standing
\ifx\paragraph\undefined\else
  \let\oldparagraph\paragraph
  \renewcommand{\paragraph}[1]{\oldparagraph{#1}\mbox{}}
\fi
\ifx\subparagraph\undefined\else
  \let\oldsubparagraph\subparagraph
  \renewcommand{\subparagraph}[1]{\oldsubparagraph{#1}\mbox{}}
\fi


\providecommand{\tightlist}{%
  \setlength{\itemsep}{0pt}\setlength{\parskip}{0pt}}\usepackage{longtable,booktabs,array}
\usepackage{calc} % for calculating minipage widths
% Correct order of tables after \paragraph or \subparagraph
\usepackage{etoolbox}
\makeatletter
\patchcmd\longtable{\par}{\if@noskipsec\mbox{}\fi\par}{}{}
\makeatother
% Allow footnotes in longtable head/foot
\IfFileExists{footnotehyper.sty}{\usepackage{footnotehyper}}{\usepackage{footnote}}
\makesavenoteenv{longtable}
\usepackage{graphicx}
\makeatletter
\def\maxwidth{\ifdim\Gin@nat@width>\linewidth\linewidth\else\Gin@nat@width\fi}
\def\maxheight{\ifdim\Gin@nat@height>\textheight\textheight\else\Gin@nat@height\fi}
\makeatother
% Scale images if necessary, so that they will not overflow the page
% margins by default, and it is still possible to overwrite the defaults
% using explicit options in \includegraphics[width, height, ...]{}
\setkeys{Gin}{width=\maxwidth,height=\maxheight,keepaspectratio}
% Set default figure placement to htbp
\makeatletter
\def\fps@figure{htbp}
\makeatother
% definitions for citeproc citations
\NewDocumentCommand\citeproctext{}{}
\NewDocumentCommand\citeproc{mm}{%
  \begingroup\def\citeproctext{#2}\cite{#1}\endgroup}
\makeatletter
 % allow citations to break across lines
 \let\@cite@ofmt\@firstofone
 % avoid brackets around text for \cite:
 \def\@biblabel#1{}
 \def\@cite#1#2{{#1\if@tempswa , #2\fi}}
\makeatother
\newlength{\cslhangindent}
\setlength{\cslhangindent}{1.5em}
\newlength{\csllabelwidth}
\setlength{\csllabelwidth}{3em}
\newenvironment{CSLReferences}[2] % #1 hanging-indent, #2 entry-spacing
 {\begin{list}{}{%
  \setlength{\itemindent}{0pt}
  \setlength{\leftmargin}{0pt}
  \setlength{\parsep}{0pt}
  % turn on hanging indent if param 1 is 1
  \ifodd #1
   \setlength{\leftmargin}{\cslhangindent}
   \setlength{\itemindent}{-1\cslhangindent}
  \fi
  % set entry spacing
  \setlength{\itemsep}{#2\baselineskip}}}
 {\end{list}}
\usepackage{calc}
\newcommand{\CSLBlock}[1]{\hfill\break\parbox[t]{\linewidth}{\strut\ignorespaces#1\strut}}
\newcommand{\CSLLeftMargin}[1]{\parbox[t]{\csllabelwidth}{\strut#1\strut}}
\newcommand{\CSLRightInline}[1]{\parbox[t]{\linewidth - \csllabelwidth}{\strut#1\strut}}
\newcommand{\CSLIndent}[1]{\hspace{\cslhangindent}#1}

\makeatletter
\@ifpackageloaded{bookmark}{}{\usepackage{bookmark}}
\makeatother
\makeatletter
\@ifpackageloaded{caption}{}{\usepackage{caption}}
\AtBeginDocument{%
\ifdefined\contentsname
  \renewcommand*\contentsname{Table of contents}
\else
  \newcommand\contentsname{Table of contents}
\fi
\ifdefined\listfigurename
  \renewcommand*\listfigurename{List of Figures}
\else
  \newcommand\listfigurename{List of Figures}
\fi
\ifdefined\listtablename
  \renewcommand*\listtablename{List of Tables}
\else
  \newcommand\listtablename{List of Tables}
\fi
\ifdefined\figurename
  \renewcommand*\figurename{Figure}
\else
  \newcommand\figurename{Figure}
\fi
\ifdefined\tablename
  \renewcommand*\tablename{Table}
\else
  \newcommand\tablename{Table}
\fi
}
\@ifpackageloaded{float}{}{\usepackage{float}}
\floatstyle{ruled}
\@ifundefined{c@chapter}{\newfloat{codelisting}{h}{lop}}{\newfloat{codelisting}{h}{lop}[chapter]}
\floatname{codelisting}{Listing}
\newcommand*\listoflistings{\listof{codelisting}{List of Listings}}
\makeatother
\makeatletter
\makeatother
\makeatletter
\@ifpackageloaded{caption}{}{\usepackage{caption}}
\@ifpackageloaded{subcaption}{}{\usepackage{subcaption}}
\makeatother
\ifLuaTeX
  \usepackage{selnolig}  % disable illegal ligatures
\fi
\usepackage{bookmark}

\IfFileExists{xurl.sty}{\usepackage{xurl}}{} % add URL line breaks if available
\urlstyle{same} % disable monospaced font for URLs
\hypersetup{
  pdftitle={Estimación de riesgo entomológico con georeferenciación en Hermosillo, Sonora, México},
  pdfauthor={Javier Edgar Verdugo Molina},
  colorlinks=true,
  linkcolor={blue},
  filecolor={Maroon},
  citecolor={Blue},
  urlcolor={Blue},
  pdfcreator={LaTeX via pandoc}}

\title{Estimación de riesgo entomológico con georeferenciación en
Hermosillo, Sonora, México}
\author{Javier Edgar Verdugo Molina}
\date{Invalid Date}

\begin{document}
\maketitle

\renewcommand*\contentsname{Table of contents}
{
\hypersetup{linkcolor=}
\setcounter{tocdepth}{2}
\tableofcontents
}
\listoffigures
\listoftables
\bookmarksetup{startatroot}

\chapter*{Resumen}\label{resumen}
\addcontentsline{toc}{chapter}{Resumen}

\markboth{Resumen}{Resumen}

Queremos Analizar los datos de la actividad de vigilancia entomológica
de forma oportuna y automática utilizando~herramientas de Ciencia de
Datos utilizando la base de datos de la plataforma de ``Vigilancia
Entomológica y Control Integral del Vector'' del INSTITUTO NACIONAL DE
SALUD PÚBLICA en colaboración con Secretaria de Salud y CENAPRECE. Con
el Fin de calcular y analizar los índices de riesgo entomológico, y con
esto gestionar información para la toma de decisiones en la prevención
de enfermedades trasmitidas por vector.

sistema integral de monitoreo de vectores
http://geosis.mx/aplicaciones/sismv/

ver {[}Garjito et al. (2020); Quintero et al. (2015); Bowman,
Runge-Ranzinger, and McCall (2014); Tun-Lin et al. (1996); Organization
(2016); South-East Asia. (2011); Tedjou et al. (2020); Braks et al.
(2003); B. Kamgang et al. (2010); Basile Kamgang et al. (2013); Weetman
et al. (2018); Tomia and Tuhatea (n.d.); Praptowibowo (2015); {]}

Vigilancia Entomológica y Control Integral del Vector
http://geosis.mx/Aplicaciones/EntomologiayControlIntegral/login.aspx

\bookmarksetup{startatroot}

\chapter{Introduccion}\label{introduccion}

\section{Motivacion}\label{motivacion}

Queremos Analizar los datos de la actividad de vigilancia entomológica
de forma oportuna y automática utilizando~herramientas de Ciencia de
Datos utilizando la base de datos de la plataforma de ``Vigilancia
Entomológica y Control Integral del Vector'' del INSTITUTO NACIONAL DE
SALUD PÚBLICA~en colaboración con Secretaria de Salud y el Centro
Nacional de Prevencion y Control de enfermedades (CENAPRECE). Con el Fin
de calcular y analizar los índices de riesgo entomológico, y con esto
gestionar información para la toma de decisiones en la prevención de
enfermedades trasmitidas por vector.

El análisis de riesgo entomológico se hace en hojas de cálculo El manejo
de la información es ineficiente: La estructura de la base de datos es
compleja Adaptar consultas a pequeñas modificaciones es impráctico La
visualización georreferenciada de esta plataforma es pobre e inestable

Usar R para visualizar la información en tiempo y espacio Automatizar el
manejo de la información Generar visualización georreferenciada con los
índices de estegomía

\section{Descripción del problema}\label{descripciuxf3n-del-problema}

El Programa estatal de vigilancia entomología y control integral de
enfermedades trasmitidas por vector, realiza diversas actividades para
la prevención de enfermedades, las cuales se guían con actividades de
vigilancia que proporcionan índices de riesgo entomológico, por ello se
hará una API para calcular los índices de riesgo entomológico de forma
automática y oportuna para la toma de decisione.

Todas las actividades se capturan diariamente en una plataforma, donde
es posible descargar los datos para su reporte y análisis. En el caso de
Vigilancia por Ovitrampa(VO), la plataforma calcula los índices de
riesgo pero para la actividad de Estudio Entomológico (EE) los calcula
por manzana trabajada, estos índices se requiere por localidad de riesgo
de manera semanal.

Se requiere poder calcular índices de riesgo del EE por tipo de estudio,
localidad de riesgo y por semana epidemiológica de manera automatizada.

La actividad de EE evalúa una muestra antes(encuesta) y
después(verificación) de las actividades de control integral para
determinar el riesgo entomológico\\
que hay en un área grande a trabajar (Sectores o Localidad) y si
requiere reforzar acciones de control y prevención.

Tener la información oportuna de los índices de riego entomológicos
permiten tomar decisiones informades de ¿En Dónde? y ¿Cuándo? focalizar
intervenciones.

\section{Gestion de datos}\label{gestion-de-datos}

Para validar y poner a prueba la solución propuesta, se plantea la
utilización de una base de datos ficticia que sea representativa con las
bases de datos reales descargadas de la plataforma. Esto permitirá
evaluar el funcionamiento del paquete de R en un entorno controlado. Una
vez realizadas las pruebas iniciales y asegurada la confiabilidad de los
resultados obtenidos,se establecerá un convenio de confidencialidad con
el fin de acceder a los datos reales y evaluar el desempeño del paquete
de R en un escenario más cercano a la realidad. Esta estrategia
garantizará la confidencialidad de los datos utilizados, a la vez que
brindará la oportunidad de comprobar la efectividad de la solución en
situaciones reales y aplicables al contexto de vigilancia entomológica y
control de enfermedades transmitidas por vectores.

\bookmarksetup{startatroot}

\chapter{Marco teorico}\label{marco-teorico}

\section{Vigilancia Epidemiológica y Entomológica del
Dengue}\label{vigilancia-epidemioluxf3gica-y-entomoluxf3gica-del-dengue}

Las enfermedades transmitidas por vectores (ETV), son padecimientos en
los que un artrópodo(mosquitos, chinches, garrapatas, etc.) actúa como
el agente infeccioso. A este tipo de artrópodos se le denomina
``vector''. Este vector actúa como hospedero y/o transmisor de la
enfermedad, mientras completa su ciclo de vida tienen diferentes
hospederos y es cuando adquiere la enfermedad y puede trasmitirla a
hospederos vertebrados como los humanos. Un ejemplo de este tipo de
enfermedades son: paludismo, dengue, leishmaniasis, oncocercosis,
tripanosomiasis, rickettsiosis, Fiebre del Oeste del Nilo, Fiebre
Chikungunya, erliquiosis, anaplasmosis y otras arbovirosis. Las ETV
representan un problema de salud pública en México. Se estima que donde
se localiza la mayor parte de centros agrícolas, ganaderos,
industriales, pesqueros, petroleros y turísticos, es decir; en cerca de
60 \% del territorio nacional, se tienen condiciones ambientales que
favorecen su transmisión (Secretaría de Salud (2023)). Entre ellas, el
dengue representa una grave amenaza, con brotes recurrentes en diversas
regiones del país. Para prevenir y controlar eficazmente las
enfermedades antes mencionadas, en México existen los programas de
vigilancia entomológica y control integral de enfermedades transmitidas
por vectores para nivel federal y estatal.

\subsection{Vigilancia
Epidemiológica:}\label{vigilancia-epidemioluxf3gica}

La vigilancia epidemiológica del dengue comprende el recopilar, analizar
y difundir sistemáticamente información sobre la enfermedad. Sus
objetivos incluyen la detección temprana de brotes, la medición de la
carga de la enfermedad, el seguimiento de las tendencias, la evaluación
del impacto social y económico. También ayuda en la asignación de
recursos, planificación de acciones preventivas y de contención. Por
ello, la vigilancia epidemiológica es esencial para el control y la
gestión eficaces contra la enfermedad.

\subsection{Índices Estegomía}\label{uxedndices-estegomuxeda}

Para evaluar la presencia y densidad de los mosquitos Aedes, los los
programas estatales de prevención y control del dengue utilizan los
llamados índices de estegomía. Estos índices cuantifican el riesgo
entomológico con respecto a criade- ros potenciales y criaderos
habitados por larvas o pupas del mosquito Aedes a través del muestreo en
campo. El muestreo debe ser uniformemente en el área de estudio y la
unidad de muestreo básica es la casa o localidad (Organization (2016)).
Al muestrear, el objetivo es observar sistemáticamente contenedores de
almacenamiento de agua. Los contenedores se examinan para detectar la
presencia de larvas y pupas de mosquitos.

La Organización Mundial de la Salud (OMS) recomienda el uso de los
índices estegomía para evaluar la presencia y la densidad de criaderos
de mosquitos en las comunidades. Los índices utilizados son el índice de
casa positiva (HI), el índice de contenedor positivo y el índice de
Breteau. En el estado de Sonora el ``Programa de Acción Específico de
Enfermedades Transmitidas por Vector''(PAE-ETV) usa estos índices para
ayudar a determinar las acciones que mitigan el riesgo entomológico. A
continuación, presento los índices de estegomía empleados por el
PAE-ETV.

\subsubsection{índice de casa positiva
(HI)}\label{uxedndice-de-casa-positiva-hi}

El índice de casa positiva (HI) es definido como el la fracción de casas
infestadas con larvas o pupas. La OMS {[}(Organization (2016)){]}
recomienda usar la siguiente expresión para su cálculo

Número de casas infestadas HI = Número de casas inspeccionadas × 100
(1.1)

En Sonora el PAE-ETV utiliza este índice HI para monitorear la
infestación de casas. Además, permiten evaluar la presencia, y estimar
la distribución del mosquito. Este índice considera a las casas como
unidad de observación, no toma en cuenta la cantidad de contenedores
positivos ni la productividad de esos contenedores.

El índice de casa positiva se ha utilizado más ampliamente para
monitorear los niveles de infestación, pero no toma en cuenta la
cantidad de contenedores positivos ni la productividad de esos
contenedores.

Dentro de una área dada, el calculo de este índice asume la misma
probabilidad de transmisión del virus. Así, si calculamos el índice HI
en una muestras de un sector se puede inferir para todo el sector. La
OMS considera en riesgo entomológico cuando al menos uno de los índices
de estegomía es igual o mayor a 5(Organization (2016))

\subsubsection{índice de contenedor positivo
(CI)}\label{uxedndice-de-contenedor-positivo-ci}

El índice de contenedor positivo (CI) se define como la proporción de
contenedores con presencia de larvas o pupas de mosquito en relación con
todos los contenedores que tienen agua (Organization, 2016). A
continuación, se muestra la siguiente expresión para su cálculo.

Número de contenedores infestados CI = Número de contenedores
inspeccionados × 100. (1.2)

El índice CI es útil para conocer la proporción de contenedores con agua
que tienen larvas o pupas de mosquito, y permite estimar el riesgo en el
área estudiada. Este índice ayuda a estimar la densidad de criaderos en
la zona. Por ello, proporciona información para identificar las áreas
con mayor riesgo de infestación. Así pues, en caso de que el índice CI
sea bajo y el HI sea alto, inferimos que el riesgo es mínimo por cada
casa muestreada, pero está muy disperso. Lo cual sugiere el realizar
acciones de control y prevención en toda el área de estudio (sector,
colonia o localidad).

El clasificar el tipo de contenedor permite elegir la estrategia
especifica para las condiciones particulares del área de estudio, como
puede ser aplicar control físico para mantener bajo control o eliminar
criaderos potenciales, dependiendo de la cantidad y el tamaño se puede
solo tirar a la basura u organizar campañas de eliminación o descacharre
d ser necesario. Otra estrategia seria el uso de químicos larvisidas
como apoyo al control fisico

\subsubsection{índice de Breteau (BI)}\label{uxedndice-de-breteau-bi}

El índice Breteau (BI) describe, número de recipientes positivos por
cada 100 casas inspeccionadas.

BI =

Número de contenedores infestados × 100 (1.3) Número de casas
inspeccionadas

El BI proporciona una relación directa entre la cantidad de recipientes
positivos y el número de casas inspeccionadas. Cunado el los
contenedores positivos tienen a ser uno por casa positiva el BI valdrá
igual al HI. En el caso que el BI salga alto pese a que los otros dos
índices sean bajos esto indica una alta concentración de recipientes por
casa aunque la proporción de casas positivas sea baja y la proporción de
recipientes positivos también sea baja, en este caso el riesgo es que
mas recipientes El BI evalúa la distribución espacial de los criaderos y
sirve para orientar las acciones de control y prevención. Sin embargo,
no refleja la productividad de los contenedores, lo que limita su
capacidad para proporcionar una evaluación completa del riesgo de
transmisión del virus. Estos datos son particularmente relevantes para
centrar esfuerzos para la gestión o eliminación de los hábitats más
comunes y para la orientación de mensajes educativos en apoyo de
iniciativas comunitarias. La vigilancia epidemiológica y entomológica
ayudan a dirigir las acciones ne- cesarias para la prevención y control
de ETVs como el dengue. Mediante la recopilación sistemática de datos
sobre la enfermedad y los vectores, se pueden tomar decisiones
informadas para mitigar su impacto y proteger a la población. los
índices de estegomía son indispensables para la vigilancia y el control
del dengue. Proporcionando información valiosa para evaluar el riesgo de
transmisión, y permiten la tomar decisiones estratégicas para priorizar
acciones y monitorear la efectividad de las intervenciones, para
prevenir la transmisión del dengue y proteger la salud de la población.

\section{Programa de Acción Específico de Enfermedades Transmitidas por
Vector}\label{programa-de-acciuxf3n-especuxedfico-de-enfermedades-transmitidas-por-vector}

El Programa de Acción Específico de Enfermedades Transmitidas por
Vector''(PAE- ETV), realiza diversas actividades para la prevención y
control de ETV. Estas actividades se pueden dividir en dos grupos: las
de vigilancia entomológicas y las de prevención y control. El objetivo
de todas estas acciones es prevenir y controlar brotes de enfermedades
transmitidas por vectores. Las ETV son aquellas causadas por un agente
vivo (como los insectos) que ingieren microorganismos de la sangre que
consumen de otros seres vivos y después los transmiten mediante sus
picaduras.

\begin{longtable}[]{@{}
  >{\raggedright\arraybackslash}p{(\columnwidth - 4\tabcolsep) * \real{0.1416}}
  >{\raggedright\arraybackslash}p{(\columnwidth - 4\tabcolsep) * \real{0.0392}}
  >{\raggedright\arraybackslash}p{(\columnwidth - 4\tabcolsep) * \real{0.8193}}@{}}
\caption{Actividades de prevencion y
control}\label{tbl-tabla-2.1}\tabularnewline
\toprule\noalign{}
\begin{minipage}[b]{\linewidth}\raggedright
Actividades
\end{minipage} & \begin{minipage}[b]{\linewidth}\raggedright
Abrebiatura
\end{minipage} & \begin{minipage}[b]{\linewidth}\raggedright
Descripcion
\end{minipage} \\
\midrule\noalign{}
\endfirsthead
\toprule\noalign{}
\begin{minipage}[b]{\linewidth}\raggedright
Actividades
\end{minipage} & \begin{minipage}[b]{\linewidth}\raggedright
Abrebiatura
\end{minipage} & \begin{minipage}[b]{\linewidth}\raggedright
Descripcion
\end{minipage} \\
\midrule\noalign{}
\endhead
\bottomrule\noalign{}
\endlastfoot
Vigilancia entomológica con ovitrampas & OV & Las ovitrampas son
contenedores con atributos ideales para que los mosquitos utilicen como
criadero así el personal del programa de vectores revisa y colecta los
huevecillos de los mosquitos, contarlos y compararlos a lo largo del año
y comparar densidad. \\
Estudios entomológicos en fase larval y pupal & EE & Búsqueda de
criaderos con larvas o pupas de mosquitos dentro de las casas de una
zona a trabajar, se realiza una búsqueda previa a las acciones de
control de ETV y una búsqueda posterior cuando el área ya ha sido
intervenida por las acciones de prevención y control. \\
Estudios Entomovirológicos & EEV & Colecta de mosquitos adultos dentro
de viviendas para identificar mosquitos transmisores de ETV, también
para el análisis por PCR de los mosquitos buscando encontrar ETV
circulando en las áreas antes de que se reporten casos en personas. \\
\end{longtable}

\begin{longtable}[]{@{}
  >{\raggedright\arraybackslash}p{(\columnwidth - 4\tabcolsep) * \real{0.1528}}
  >{\raggedright\arraybackslash}p{(\columnwidth - 4\tabcolsep) * \real{0.1528}}
  >{\raggedright\arraybackslash}p{(\columnwidth - 4\tabcolsep) * \real{0.6944}}@{}}
\caption{Actividades de prevencion y
control}\label{tbl-tabla-2.2}\tabularnewline
\toprule\noalign{}
\begin{minipage}[b]{\linewidth}\raggedright
Actividades
\end{minipage} & \begin{minipage}[b]{\linewidth}\raggedright
Abrebiatura
\end{minipage} & \begin{minipage}[b]{\linewidth}\raggedright
Descripcion
\end{minipage} \\
\midrule\noalign{}
\endfirsthead
\toprule\noalign{}
\begin{minipage}[b]{\linewidth}\raggedright
Actividades
\end{minipage} & \begin{minipage}[b]{\linewidth}\raggedright
Abrebiatura
\end{minipage} & \begin{minipage}[b]{\linewidth}\raggedright
Descripcion
\end{minipage} \\
\midrule\noalign{}
\endhead
\bottomrule\noalign{}
\endlastfoot
Control Larvario & CL & Búsqueda intencionada de criaderos potenciales
de mosquitos de un área determinada para la aplicación de medidas de
prevención, en todas las casas que permitan la entrada. \\
Rociado a caso probable & RaC & Aplicación intradomiciliar de
insecticida con efecto residual en casas dentro de un cerco de 5 a 9
casas en respuesta a un caso probable o confirmado a Dengue, Zika o
Chicungunya. \\
Rociado intradomiciliario & RI & Aplicación intradomiciliar de
insecticida con efecto residual en casas de un área que se considere de
riesgo entomológico y susceptible a contagios de ETV. \\
Nebulización espacial & N & Aplicación de rociado espacial de
insecticidas sin efecto residual en un área con alto riesgo entomológico
o en respuesta a casos probables o confirmados. \\
\end{longtable}

\subsection{Plataforma de Vigilancia Entomológica y Control Integral del
Vector}\label{plataforma-de-vigilancia-entomoluxf3gica-y-control-integral-del-vector}

El Programa de Acción Específico de Enfermedades Transmitidas por Vector
(PAE-ETV) registra diariamente en una plataforma todas las actividades
men- cionadas en las tablas Tabla 1.1 y Tabla 1.2. Esto permite al
programa estatal descargar y analizar los datos. Para la actividad de
EE, el programa calcula los resultados por área trabajada y requiere el
reporte semanal con base en la ubicación, como se observa en la Tabla
2.2 y el sistema integral de vigilancia de vectores (2023). La actividad
de EE implica el muestreo de formas inmaduras de los mosquitos (las
larvas y las pupas) en lugar de capturar huevos o mosqui- tos adultos.
Se realizan búsquedas sistemáticas para detectar recipientes que
contengan agua e inspeccionar la presencia de larvas, pupas y larvas y
pupas de mosquitos. Esta actividad evalúa una muestra antes (encuesta) y
después (verificación) de las actividades de control integral para
determinar el riesgo entomológico en un área trabajada (sectores o
localidades) y si es necesario reforzar las acciones de control, ver
Tabla 1.1 y Tabla 1.2. riesgo {[}Sistema integral de monitoreo de
vectores (2023);{]}

\subsection{Datos Geograficos}\label{datos-geograficos}

El PAE-ETV proporcionó shapefiles del estado de Sonora y como esta todo
dividido en sectores. Estos son los sectores en los que se basa la
distribución de todas las actividades del PAE-ETV y se usan para poder
geopreferencia las áreas intervenidas, determinar el riesgo entomológico
y planear las actividades de control y prevención que se requieran. Los
shapefiles pueden usarse para crear mapas con los indines de estegomía
específicos en tiempo y lugar según se necesite.

\section{Desarrollo de un paquete de R con desarrollo basado en pruebas
(TDD)}\label{desarrollo-de-un-paquete-de-r-con-desarrollo-basado-en-pruebas-tdd}

Para determinar el riesgo entomológico delimitado por área geográfica
(desde un sector específico o hasta una localidad o municipio).Es
necesario obtener la información de la vigilancia entomología de forma
periódica, y comprender la estructurada la base de datos de EE. Por lo
cual implementaremos herramientas de ciencias de datos para aplicar una
metodología en la gestión y análisis de esta información. Esto
implicaría calcular los índices de estegomia para un lugar y en un
tiempo determinado para la toma de decisiones. Este trabajo presenta el
desarrollo de un paquete R para el análisis de los datos de vigilancia
entomológica del dengue, siguiendo una metodología de desarrollo basado
en pruebas (TDD). El paquete combina la potencia de Tidyverse para la
manipulación de datos, Devtools para la gestión de Testthat para la
escritura de pruebas unitarias

El desarrollo de software en R ha evolucionado hacia metodologías más
sofistica- das y automatizadas. El uso combinado de frameworks como
Tidyverse, Devtools y Testthat permite crear un trabajo con desarrollo
basado en pruebas (TDD) para proyectos de R (ver Figura 1.1). Tidyverse
es una colección de paquetes R enfocados en la manipulación de datos,
análisis estadístico y visualización. Su filosofía funcional facilita la
escritura de código modular y limpio. Devtools es un paquete R
indispensable para la gestión de proyectos, creación de paquetes y
publicación en CRAN. Automatiza tareas comunes como la ejecución de
pruebas y generación de documentación. Testthat es un paquete R crucial
para implementar TDD. Permite escribir pruebas unitarias que garantizan
el funcionamiento correcto de cada función o módulo del código.

El desarrollo basado en pruebas (TDD) es una metodología de desarrollo
de software en la que se escriben pruebas antes de escribir el código.
Esto ayuda a garantizar que el código esté bien escrito y cumpla con los
requisitos. Escribir pruebas: Se definen pruebas unitarias para cada
función utilizando Testthat. Desarrollar código: Se implementa el código
que satisface las pruebas escritas. Ejecutar pruebas: Se ejecutan las
pruebas para verificar el funciona- miento del código. Refactorizar: Se
mejora la legibilidad y mantenibilidad del código sin alterar su
funcionalidad. Repetir: Se repiten los pasos anteriores hasta completar
el desarrollo. Combinar Tidyverse, Devtools y Testthat en un flujo de
trabajo TDD es una práctica recomendada para el desarrollo de software
en R. Esta metodología per- mite tener menos errores, crear código más
robusto, legible y eficiente, facilitando la colaboración y el
mantenimiento a largo plazo de los proyectos. El paquete R presentado en
este trabajo constituye una valiosa herramienta para la vigilancia
entomológica del dengue. Su uso puede contribuir a mejorar la toma de
decisiones en materia de prevención y control del dengue, fortaleciendo
las estrategias de salud pública en México.

\bookmarksetup{startatroot}

\chapter{Prelimilares}\label{prelimilares}

\section{descripcion del
problematica}\label{descripcion-del-problematica}

\section{deacuerdo al marco teorico como se
soliciona}\label{deacuerdo-al-marco-teorico-como-se-soliciona}

\section{}\label{section}

\bookmarksetup{startatroot}

\chapter{Cálculo de los índices de
estegomía}\label{cuxe1lculo-de-los-uxedndices-de-estegomuxeda}

Índice de casa (local) (IC), es decir, porcentaje de casas infestadas
con larvas, pupas o ambas. CI = Casas infestadas/casas inspeccionadas X
100

Índice de recipiente (IR), es decir, porcentaje de recipientes que
contienen agua y están infestados con larvas o pupas. IR = Recipientes
positivos / recipientes inspeccionados X 100

Dengue: guías para el diagnóstico, tratamiento, prevención y control.
Índice de Breteau (IB) -- es decir, número de recipientes positivos por
cada 100 casas inspeccionadas. IB = Número de recipientes positivos /
casas inspeccionadas X 100

\section{}\label{section-1}

Se sigue el modelo de desarrollo basado test driven development Se
desarrollaron basado en tidyr y devtools tidyr Manipulación de los datos
Estilo de codificación devtools Documentación (roxygen2 ) Pruebas
unitarias (testthat) CRAN (Comprehensive R Archive Network)

\bookmarksetup{startatroot}

\chapter{Limpieza de datos}\label{limpieza-de-datos}

Los datos para calcular índices de riesgo entomológico son de las
actividades de vigilancia entomológica con ovitrampas y estudios
entomológicos para fase larvaria y pupal. Estos datos se descargan de la
plataforma en archivos .txt y requieren de un tratamiento previo, que
incluye un código particular para leer el tipo de archivo, selección de
variables de interés, cambiar el formato de los títulos de las variables
y cambios en los tipos de datos, antes de poder realizar cálculos y
análisis.

La base de datos de Estudio entomológico se descarga en un archivo .txt
de la Plataforma ``Vigilancia Entomológica y Control Integral del
Vector'', tiene 121 variables de las cuales consideramos 16 de interés,
3 de ellas se dividen en dos por lo que resultan en 19

\section{descripcion de la base de
datos}\label{descripcion-de-la-base-de-datos}

La base de datos de estudio entomológico es un archivo .txt de la
Plataforma Vigilancia Entomológica y Control Integral del Vector

Tiene 121 variables

y 16 variables de interés

formato de los datos

\section{variables de interes}\label{variables-de-interes}

\subsection{Variable para determinar el tipo de estudio
entomológico:}\label{variable-para-determinar-el-tipo-de-estudio-entomoluxf3gico}

Tipo de Estudio

\subsection{Variables para la
georreferenciación:}\label{variables-para-la-georreferenciaciuxf3n}

Jurisdicción Municipio Localidad Sector

\subsection{Variables de intervalos de
tiempo:}\label{variables-de-intervalos-de-tiempo}

Fecha de Inicio Semana Epidemiológica

\subsection{Variables para índices de
estegomía:}\label{variables-para-uxedndices-de-estegomuxeda}

Casas Revisadas Casas Positivas Total de Recipientes con Agua Total de
Recipientes Positivos Total de Recipientes Positivos a Pupas No.~Total
de Pupas en Recipientes Recipientes Tratables Recipientes Controlables
Recipientes Eliminables

\section{Archivos shapefile de los sectores de todo el estado de
Sonora}\label{archivos-shapefile-de-los-sectores-de-todo-el-estado-de-sonora}

\section{Data cleaning}\label{data-cleaning}

Resolución de inconsistencia de tipo de dato:

Formato de títulos de las variables Tipo de dato POSIX~a las variables
de fecha Asignar tipo de dato correcto(carácter) Dividir variables
compuestas

Ajuste entre tipos de datos shapefile para empatar con un dataframe
Carga el archivo .shp (sp) Cambiar la proyección EPSG 4326(sp) Guardar
como un archivo .shp\\
Carga el archivo .shp (maptools) Transforma en un dataframe con datos
geoespaciales

````descripcion de la base de datos formato de los datos variables de
interes Data cleaning''''

\bookmarksetup{startatroot}

\chapter{Referencias}\label{referencias}

\phantomsection\label{refs}
\begin{CSLReferences}{1}{0}
\bibitem[\citeproctext]{ref-Bowman2014}
Bowman, Leigh R., Silvia Runge-Ranzinger, and P. J. McCall. 2014.
{``Assessing the Relationship Between Vector Indices and Dengue
Transmission: A Systematic Review of the Evidence.''} \emph{PLoS
Neglected Tropical Diseases} 8.
\url{https://doi.org/10.1371/JOURNAL.PNTD.0002848}.

\bibitem[\citeproctext]{ref-Braks2003}
Braks, Marieta A. H., Nildimar A. Honório, Ricardo Lourenço-De-Oliveira,
Steven A. Juliano, and L. Philip Lounibos. 2003. {``Convergent Habitat
Segregation of Aedes Aegypti and Aedes Albopictus (Diptera: Culicidae)
in Southeastern Brazil and Florida.''} \emph{Journal of Medical
Entomology} 40: 785--94.
\url{https://doi.org/10.1603/0022-2585-40.6.785}.

\bibitem[\citeproctext]{ref-Garjito2020}
Garjito, Triwibowo Ambar, Muhammad Choirul Hidajat, Revi Rosavika
Kinansi, Riyani Setyaningsih, Yusnita Mirna Anggraeni, Mujiyanto, Wiwik
Trapsilowati, et al. 2020. {``Stegomyia Indices and Risk of Dengue
Transmission: A Lack of Correlation.''} \emph{Frontiers in Public
Health} 8 (July). \url{https://doi.org/10.3389/fpubh.2020.00328}.

\bibitem[\citeproctext]{ref-Kamgang2013}
Kamgang, Basile, Carine Ngoagouni, Alexandre Manirakiza, Emmanuel
Nakouné, Christophe Paupy, and Mirdad Kazanji. 2013. {``Temporal
Patterns of Abundance of Aedes Aegypti and Aedes Albopictus (Diptera:
Culicidae) and Mitochondrial DNA Analysis of Ae. Albopictus in the
Central African Republic.''} \emph{PLoS Neglected Tropical Diseases} 7.
\url{https://doi.org/10.1371/journal.pntd.0002590}.

\bibitem[\citeproctext]{ref-Kamgang2010}
Kamgang, B., J. Y. Happi, P. Boisier, F. Njiokou, J. P. Hervé, F.
Simard, and C. Paupy. 2010. {``Geographic and Ecological Distribution of
the Dengue and Chikungunya Virus Vectors Aedes Aegypti and Aedes
Albopictus in Three Major Cameroonian Towns.''} \emph{Medical and
Veterinary Entomology} 24 (June): 132--41.
\url{https://doi.org/10.1111/j.1365-2915.2010.00869.x}.

\bibitem[\citeproctext]{ref-Organization2016}
Organization, World Health. 2016. \emph{Entomological Surveillance for
Aedes Spp. In the Context of Zika Virus. Interim Guidance for
Entomologists}. WHO Regional Office for South-East Asia.
\url{https://iris.who.int/handle/10665/204894}.

\bibitem[\citeproctext]{ref-Praptowibowo2015}
Praptowibowo, Wahyu. 2015. {``Maya Index Dan Gambaran Habitat
Perkembangbiakan Larva Aedes Sp. Berdasarkan Endemisitas DBD Di Kota
Semarang Provinsi Jawa Tengah''} 3: 2356--3346.
\url{http://ejournal-s1.undip.ac.id/index.php/jkm}.

\bibitem[\citeproctext]{ref-Quintero2015}
Quintero, J., T. Garcia-Betancourt, S. Cortes, D. Garcia, L. Alcala, C.
Gonzalez-Uribe, H. Brochero, and G. Carrasquilla. 2015. {``Effectiveness
and Feasibility of Long-Lasting Insecticide-Treated Curtains and Water
Container Covers for Dengue Vector Control in Colombia: A Cluster
Randomised Trial.''} \emph{Transactions of the Royal Society of Tropical
Medicine and Hygiene} 109 (February): 116--25.
\url{https://doi.org/10.1093/trstmh/tru208}.

\bibitem[\citeproctext]{ref-SouthEastAsia.2011}
South-East Asia., World Health Organization. Regional Office for. 2011.
\emph{Comprehensive Guidelines for Prevention and Control of Dengue and
Dengue Haemorrhagic Fever.} World Health Organization Regional Office
for South-East Asia. \url{https://iris.who.int/handle/10665/204894}.

\bibitem[\citeproctext]{ref-Tedjou2020}
Tedjou, Armel N., Basile Kamgang, Aurélie P. Yougang, Theodel A.
Wilson-Bahun, Flobert Njiokou, and Charles S. Wondji. 2020. {``Patterns
of Ecological Adaptation of Aedes Aegypti and Aedes Albopictus and
Stegomyia Indices Highlight the Potential Risk of Arbovirus Transmission
in Yaoundé, the Capital City of Cameroon.''} \emph{Pathogens} 9 (June):
491. \url{https://doi.org/10.3390/pathogens9060491}.

\bibitem[\citeproctext]{ref-Tomia}
Tomia, Amalan, and Rosmila Tuhatea. n.d. {``Potential of Transmission of
the Dengue Virus Based on Entomological Index and Maya Index in Kalumata
and North Mangga Dua Villages,ternate City.''} \emph{International
Journal of Science}. \url{http://ijstm.inarah.co.id}.

\bibitem[\citeproctext]{ref-TunLin1996}
Tun-Lin, W., B. H. Kay, A. Barnes, and S. Forsyth. 1996. {``Critical
Examination of Aedes Aegypti Indices: Correlations with Abundance.''}
\emph{American Journal of Tropical Medicine and Hygiene} 54.
\url{https://doi.org/10.4269/ajtmh.1996.54.543}.

\bibitem[\citeproctext]{ref-Weetman2018}
Weetman, David, Basile Kamgang, Athanase Badolo, Catherine L. Moyes,
Freya M. Shearer, Mamadou Coulibaly, João Pinto, Louis Lambrechts, and
Philip J. McCall. 2018. {``Aedes Mosquitoes and Aedes-Borne Arboviruses
in Africa: Current and Future Threats.''} \emph{International Journal of
Environmental Research and Public Health}. MDPI.
\url{https://doi.org/10.3390/ijerph15020220}.

\end{CSLReferences}



\end{document}
